\section{The \textsc{RippleEval} Benchmark}
\label{sec:dataset}

\yuji{Just updated a new version of the benchmark, will fix more details in more rounds later}

To comprehensively analyze how factual updates affect related knowledge after post-training, we construct \textbf{\textsc{RippleEval}}, a QA-based benchmark grounded in verified factual relations. 

Each QA sample in \textsc{RippleEval} corresponds to an underlying factual relation $(s,r,o)$, where $s$ and $o$ are entities and $r$ is a relation. We use these grounded relations to connect QA samples because QA expressions themselves are not stable units for defining factual dependency: the same fact can be asked in different ways, while unrelated facts may share similar question templates. Thus, the graph is used to define factual connections among QA samples, while model training and evaluation are still performed through natural-language QA.

\paragraph{Grounded Graph Construction.}
We construct an entity-level factual graph and use it to organize QA samples. In this graph, each node is an entity, and each directed edge represents a verified factual relation $(s,r,o)$ from subject entity $s$ to object entity $o$ under relation $r$. Each verified edge is then verbalized into a QA sample for training and evaluation.


\paragraph{Relation Constraints.}
To make factual connections and downstream effects well-defined, we restrict \textsc{RippleEval} to relations with a unique or primary expected answer. Existing benchmarks such as MQUAKE \cite{zhong2023mquake} and RippleEdits \cite{cohen2023ripple} evaluate whether edited facts lead to correct multi-hop or ripple-effect consequences. However, their goal is broad evaluation of editing consistency rather than controlling whether each subject-relation pair has a single expected object answer. In \textsc{RippleEval}, we impose this constraint to make the affected factual neighbors and their expected answers unambiguous.

Specifically, for each subject-relation pair $(s,r)$, we require one expected object answer $o$. Relations such as \textit{CapitalOf} and \textit{DevelopedByPrimary} satisfy this requirement, while relations such as \textit{HasChild} are excluded because they can have many valid objects. Note that it only requires each subject-relation pair to have a single expected answer. This constraint allows us to define controlled factual neighborhoods for measuring ripple effects.

\paragraph{Expansion and Verification.}
We employ an automated pipeline to construct the factual grounding graph, $\mathcal{G}_{fact}$. We start from high-confidence seed triplets, such as \texttt{("Minecraft", "DevelopedByPrimary", "Mojang Studios")}, and expand the graph by following selected factual relations from the current entities. During expansion, we avoid cycles so that a downstream sample does not point back to an earlier sample in the same path. This keeps the direction of factual neighborhoods clear when we later evaluate ripple effects. 

\paragraph{Scaling Up Knowledge Graph.}
Instead of relying on small-scale local models, we utilize the high-throughput DeepSeek API (\texttt{deepseek-chat}) to simultaneously propose candidate triplets and generate high-quality natural language questions. Scaling this rigorous generation process consumed approximately 309.2 million tokens (182.6M prompt, 126.6M completion). Crucially, the generated questions are hardcoded directly into the graph's edge attributes. This mechanism strictly eliminates prompt variance during evaluation: models are tested using the exact same question string before and after the update, ensuring failures stem from structural corruption rather than wording ambiguity. Each candidate is then verified against Wikidata \cite{vrandecic2014wikidata} through an external validation module, and unverified triplets are discarded. The final graph $\mathcal{G}_{fact}$ comprises $100,015$ unique entity nodes and $432,562$ verified relation edges spanning 39 strict relation types.

\paragraph{Factual Connectivity as Popularity.}
After constructing the verified factual graph, we use entity connectivity to approximate the popularity of factual knowledge in the benchmark. For a factual relation $(s,r,o)$, we measure the connectivity of its object entity $o$ by its in-degree, i.e., how many verified subject-relation facts point to the same object. This score does not measure how often a QA wording appears; instead, it measures how widely the answer entity is shared across different factual relations.

This definition gives each QA sample a grounded popularity score through its underlying factual relation. A sample whose answer entity is connected to many other facts is treated as involving a more widely shared factual anchor, while a sample whose answer entity has few incoming facts is treated as involving a narrower factual context. In this way, the graph connectivity serves as a proxy for the uneven distribution of real-world knowledge, where some entities are linked to many facts and others are much more isolated.

